\section{Divisibility at inert primes}\label{sec:divisibility}

The gains of \cref{sec:cast} are paid for by one phenomenon: Taylor coefficients of theta
functions at a CM point are highly divisible at primes inert in the CM field. The note states
it as: a coefficient of order $a$, divided by $\Omega^a$, is divisible by
$p^{a/p+a/p^3+\cdots-O(\log a)}$. In this section $\ell$ is such a prime and
\[
\gamma_\ell=\frac{\ell}{\ell^2-1}=\frac1\ell+\frac1{\ell^3}+\frac1{\ell^5}+\cdots,\qquad
\rho_\ell=\frac1{\ell^2-1}.
\]

\subsection{The mechanism}
Let $F$ be the formal group over $\Z_\ell$ whose logarithm is \lvS{} (\lean{ltLog})
\[
\log_F(X)=\sum_{j\ge0}\frac{X^{\ell^{2j}}}{\ell^{j}},
\]
a Lubin--Tate formal group of height two \cite{lubintate1965,honda1970}. Its exponential has
$r$-th coefficient of valuation at least $-\rho_\ell(r-1)$ \lvS{} (\lean{LemmaSt}). Hence if
$\vartheta$ is a power series with $\ell$-integral coefficients in $X$ and we expand it in the
variable $z=\upsilon\log_F(X)$ with $\upsilon$ a unit, the $r$-th Taylor coefficient of
$\vartheta$ in $z$, multiplied by $r!$, has valuation at least
\[
v_\ell(r!)-\rho_\ell\,r\ \ge\ \frac{r}{\ell-1}-\frac{r}{\ell^2-1}-(\log_\ell r+1)
\ =\ \gamma_\ell\,r-(\log_\ell r+1).
\]
That is the whole reason for the exponent $a/\ell+a/\ell^3+\cdots$: derivatives in the
logarithmic coordinate of a height-two formal group gain a factor $\ell^{\gamma_\ell}$ per
order. The archive's general form \lvS{} (\lean{LemmaRA}) bounds, for a function
$U(\mathsf g(z))\exp(\pm\mathsf b(z))$
in several variables, the valuation of a derivative of multi-order $\kappa$ by
$\sum_i\kappa_i(\gamma_\ell-\varsigma_i)$ less a logarithmic loss, with weights
$\varsigma_i\ge0$, and a second theorem treats
several ``letters'' at once, where the top terms take the form $\sum_{e\text{ odd}}\nu/\ell^e$
cut off by a capacity: the same shape as the count \eqref{eq:ask}. Both are on the path of the
final theorem \lvM. The exponential factor is of the kind that Dwork introduced
\cite{dwork1960}.

\subsection{Why the function is integral: complex multiplication}
A CM elliptic curve really is used; the argument is not carried out with power series alone.
\begin{statement}[\lean{CM2}]\label{st:cm2}
Let $d$ be squarefree with $d\equiv3\pmod4$ and let $L$ be a lattice with complex
multiplication by $\Z[(1+\sqrt{-d})/2]$. For every prime $\ell\ge5$ not dividing $d$ and every
isomorphism $\iota\colon\C\cong\C_\ell$ there is $\xi\ne0$ such that
$\xi^4\iota g_2(L)$ and $\xi^6\iota g_3(L)$ are integral in the maximal unramified
extension of $\Q_\ell$ and the discriminant of the resulting Weierstrass equation is a unit.
\end{statement}
\begin{statement}[\lean{CM3}]\label{st:cm3}
In the situation of \cref{st:cm2}, suppose moreover that $-d$ is not a square modulo $\ell$.
Then there is a unit $\upsilon$ of the maximal unramified extension of $\Q_\ell$ such that the
local series of the curve, composed with $\upsilon\log_F$, has integral coefficients in that
extension.
\end{statement}
\Cref{st:cm2} is good reduction \cite{serretate1968}; the archive derives it from global Artin
reciprocity \lvS{} (\lean{cmGoodModel_of_globalArtinReciprocity}), which it also proves.
\Cref{st:cm3} says that at an inert prime the formal group of the curve is the height-two group
$F$, so that the complex uniformising coordinate $z$ is a unit times $\log_F$ of an integral
parameter. The reduction is supersingular by Deuring's criterion \cite{deuring1941}, and a
formal group of height two with this action of the ring of integers is the Lubin--Tate group
$F$ up to isomorphism \cite{lubintate1965,honda1970}.

From these the archive proves, for every prime $\ell\ge5$ with $\bigl(\tfrac{-d}{\ell}\bigr)=-1$ and
every embedding of $\overline\Q$ into $\C_\ell$, that the kernel of the theta lift attached to a
pair of admitted lattices is integral \lvS{} (\lean{IntInert}). A \emph{jet} is the value at
the origin of a differential operator applied to the Taylor series of the kernel; the periods
of \cref{sec:adm} are expressed through jets, and that is how the divisibility reaches them.

The jet theorem itself is applied to an arbitrary series satisfying the integrality hypothesis
\lvP{} (\lean{prvF_of_orbitLetterBounds}); the work is in showing that the actual kernel
satisfies it. The letters of the general theorem index the positions along an orbit of
Frobenius at $\ell$ acting on the embeddings of the CM field.

\subsection{Integrality of the periods}
The periods themselves, for forms on the definite group, are finite sums over the classes of a
genus. Their integrality rests on a statement about representations of general linear groups \lvS{}
(\lean{standard_weyl_lattices_fock_dual}): the integral forms of the Weyl module and of the
dual Weyl module are exact duals of each other under the pairing of the polynomial (Fock)
model; compare \cite{carterlusztig1974}. Denominators coming from orders of finite groups are
controlled by Minkowski's theorem \cite{minkowski1887}, in the form \lvS{}
(\lean{minkowski_bound_rat}): the order of a finite subgroup of $\GL_N(\Q)$ divides an
explicit integer depending on $N$ alone.
